\section{The hollow-core Hamiltonian}\label{sec:hollowcore}

\subsection{The main theorem}

\begin{theorem}[no uniform gap above the zero modes]\label{thm:main}
Fix $\lambda>0$. For even $N$ let $\Tor_N$ be the rectangular torus with $L_x/L_y=\lambda$ and $L_xL_y=4\pi N$, let
$H_N$ be the hollow-core Hamiltonian of Definition~\ref{def:hamiltonian}, and let $E_1(N)\le E_2(N)\le\cdots$ be the
eigenvalues of $H_N$ in the sector $\sect{(N/2,N/2)}$, repeated according to multiplicity, so that
$E_1(N)=\cdots=E_{10}(N)=0<E_{11}(N)$.
\begin{enumerate}[(i)]
\item If for some $p\in\N$, $\Delta\ge0$ and $\gamma>2\Delta$ the inequalities $E_p(N)-E_1(N)\le\Delta$ and
$E_{p+1}(N)-E_p(N)\ge\gamma$ hold for infinitely many even $N$, then $p\in4\Z$ and $p\ge12$.
\item $E_{11}(N)\to0$ as $N\to\infty$ through the even integers.
\end{enumerate}
\end{theorem}

\begin{proof}
The side lengths are $L_x=(4\pi\lambda N)^{1/2}$ and $L_y=(4\pi N/\lambda)^{1/2}$. Both tend to infinity with $N$, so
there is $N_\lambda$ such that the torus is admissible, that is $L_y\ge L_0$ and $L_x\ge4\delta_0+4$
(Definition~\ref{def:frame}), for every $N\ge N_\lambda$. Put $c_*:=(4\pi\lambda)^{1/2}$, so that $L_x=c_*\sqrt N$.

(i) Let $\mathcal N$ be the set of even $N\ge N_\lambda$ for which the two inequalities hold; it is infinite. The tori
$(\Tor_N)_{N\in\mathcal N}$ are admissible with $L_x\ge c_*\sqrt N$, and the hollow-core Hamiltonians
$(H_N)_{N\in\mathcal N}$ satisfy (H1) and (H2) by Corollary~\ref{cor:hc-in-class}; so Setting~\ref{set:index} holds. The two
inequalities say that Assumption~\ref{ass:gap} holds for the sectors $\nn_N=(N/2,N/2)$, $N\in\mathcal N$, with $p$,
$\Delta$ and $\gamma$ independent of $N$. Corollary~\ref{cor:classwide} gives $p\in4\Z$. If $p\le9$, then
$E_p(N)=E_{p+1}(N)=0$ by \eqref{eq:E11}, which contradicts $E_{p+1}(N)-E_p(N)\ge\gamma>0$. Hence $p\ge10$, and since
neither $10$ nor $11$ lies in $4\Z$, $p\ge12$.

(ii) Suppose not. Then there are $\gamma>0$ and infinitely many even $N$ with $E_{11}(N)\ge\gamma$. By \eqref{eq:E11},
for these $N$ we have $E_{10}(N)-E_1(N)=0$ and $E_{11}(N)-E_{10}(N)=E_{11}(N)\ge\gamma$, so the hypothesis of (i) holds with
$p=10$, $\Delta=0$ and this $\gamma$. By (i), $10\in4\Z$, which is false.
\end{proof}

The kernel enters the proof only through its dimension, ten, and through $E_1(N)=0$ (Theorem~\ref{thm:kernel}); no
other property of the Haldane--Rezayi states is used. In the language of Section~\ref{sec:index}, along a sequence as
in the proof of (ii) the transported charge is $\expP{T_-}=-\expP{n_{0\up}}=-\frac14$ (Lemma~\ref{lem:theta}(d)), and
$10\cdot\frac14=\frac52$ is at distance $\frac12$ from $\Z$.

\begin{corollary}[odd $N$]\label{cor:oddN}
Fix $\lambda>0$, and for odd $N$ let $\Tor_N$ and $H_N$ be as in Theorem~\ref{thm:main}. In each of the sectors
$\sect{((N\pm1)/2,(N\mp1)/2)}$ let $E_1(N)\le E_2(N)\le\cdots$ be the eigenvalues of $H_N$, so that
$E_1(N)=E_2(N)=0<E_3(N)$.
\begin{enumerate}[(i)]
\item $E_3(N)\to0$ as $N\to\infty$ through the odd integers.
\item For every $D\in\N$ and $\gamma>0$, for all sufficiently large odd $N$ there are more than $D$ eigenvalues in
$[0,\gamma)$.
\end{enumerate}
\end{corollary}

\begin{proof}
Let $N_\lambda$ and $c_*$ be as in the proof of Theorem~\ref{thm:main}, and let $\mathcal N$ be the set of odd
$N\ge N_\lambda$. As there, the tori $(\Tor_N)_{N\in\mathcal N}$ and the hollow-core Hamiltonians are as in
Setting~\ref{set:index}, with (H1) and (H2). We treat the sectors $\nn^+_N=((N+1)/2,(N-1)/2)$; the sectors $\nn^-_N$ are
treated in the same way. By Theorem~\ref{thm:kernel}(ii) and $H_N\ge0$, the eigenvalues of $H_N$ in $\sect{\nn^+_N}$
satisfy $E_1(N)=E_2(N)=0<E_3(N)$.

(i) Suppose that there are $\gamma>0$ and infinitely many odd $N$ with $E_3(N)\ge\gamma$. For these $N$,
$E_2(N)-E_1(N)=0$ and $E_3(N)-E_2(N)=E_3(N)\ge\gamma$, so Assumption~\ref{ass:gap} holds for the sectors $\nn^+_N$ along the
infinite set of such $N$ with $N\ge N_\lambda$, with $p=2$, $\Delta=0$ and this $\gamma$. This contradicts
Corollary~\ref{cor:classwide-odd}(i).

(ii) By Corollary~\ref{cor:classwide-odd}(ii), for every $D\in\N$ and $\gamma>0$ and all but finitely many
$N\in\mathcal N$, the restriction of $H_N$ to $\sect{\nn^+_N}$ has more than $D$ eigenvalues, counted with multiplicity,
in $[E_1(N),E_1(N)+\gamma)=[0,\gamma)$.
\end{proof}

By Remark~\ref{rem:class-scope}, both statements reflect the unpaired spin forced by the density of these sectors;
they say nothing about the sector $(N/2,N/2)$ at even $N$.

\subsection{The scope over aspect ratios}

The construction of Section~\ref{sec:frame} distinguishes the two sides of the torus: the Landau-gauge orbitals are
localized across the side of length $L_x$, and admissibility asks $L_y\ge L_0$ and $L_x\ge4\delta_0+4$. The spectrum of
the hollow-core Hamiltonian does not distinguish them.

\begin{lemma}[exchange of the sides]\label{lem:exchange}
Let $L_x,L_y>0$ with $L_xL_y=4\pi N$, and let $\Tor'_N$ be the torus of Definition~\ref{def:torus} with side lengths
$L'_x:=L_y$ and $L'_y:=L_x$. Denote its lowest Landau level, its Fock space, the annihilation operators on it, its number
operators and its hollow-core Hamiltonian by $\mathfrak h'_N$, $\mathcal H'_N$, $c'(\cdot)$, $N'_\sigma$ and $H'_N$. For
$f\in\hN$ let $(\mathcal Rf)(x,y):=e^{\iu xy}f(-y,x)$. Then $\mathcal R$ is a unitary from $\hN$ onto $\mathfrak h'_N$; the
unitary $\Ff(\mathcal R\otimes\one):\Fock\to\mathcal H'_N$ that fixes the vacuum and satisfies
$\Ff(\mathcal R\otimes\one)\,c(f)\,\Ff(\mathcal R\otimes\one)^*=c'((\mathcal R\otimes\one)f)$ for $f\in\kN$ maps $N_\sigma$ to
$N'_\sigma$, $\sigma\in\{\up,\dn\}$, and $H_N$ to $H'_N$. In particular, for every sector label $\nn$, the restrictions of
$H_N$ to $\sect\nn$ and of $H'_N$ to the sector of $\mathcal H'_N$ with the same label have the same eigenvalues, with
multiplicities.
\end{lemma}

\begin{proof}
\emph{The one-particle map.} Let $\varrho(x,y):=(-y,x)$, a rotation of $\R^2$. It maps the lattice $L'_x\Z\times L'_y\Z$
onto $L_x\Z\times L_y\Z$, since $\varrho(mL_y,nL_x)=(-nL_x,mL_y)$, and hence induces a measure-preserving bijection of
$\Tor'_N$ onto $\Tor_N$. Consequently, for every smooth function $F$ on $\Tor_N$, the function $F\circ\varrho$ is smooth on
$\Tor'_N$, $\int_{\Tor'_N}F\circ\varrho\,\dd r=\int_{\Tor_N}F\,\dd r$, and $\nabla^2(F\circ\varrho)=(\nabla^2F)\circ\varrho$, because
the Laplacian commutes with rotations.

Let $f\in\hN$ and $g:=\mathcal Rf$; $g$ is smooth. By the two conditions \eqref{eq:bc} for $f$,
\begin{gather*}
g(x,y+L'_y)=e^{\iu x(y+L_x)}f(-y-L_x,x)=e^{\iu xy}e^{\iu L_xx}e^{-\iu L_xx}f(-y,x)=g(x,y),\\
g(x+L'_x,y)=e^{\iu(x+L_y)y}f(-y,x+L_y)=e^{\iu L'_xy}g(x,y),
\end{gather*}
which are the conditions \eqref{eq:bc} for $\Tor'_N$. With $D_x=\partial_x$ and $D_y=\partial_y-\iu x$, and with
$\partial_1f$, $\partial_2f$ the partial derivatives of $f$ in its two arguments,
\begin{align*}
e^{-\iu xy}\bigl((D_x+\iu D_y)g\bigr)(x,y)&=\iu y\,f(-y,x)+(\partial_2f)(-y,x)-\iu(\partial_1f)(-y,x)\\
&=-\iu\bigl(\partial_1f+\iu\partial_2f+s\,f\bigr)(s,t)\Big|_{(s,t)=(-y,x)}=-\iu\bigl((D_x+\iu D_y)f\bigr)(-y,x)=0,
\end{align*}
by \eqref{eq:lll} for $f$. So $g\in\mathfrak h'_N$. Moreover $\overline{\mathcal Rf_1}\,\mathcal Rf_2=(\overline{f_1}f_2)\circ\varrho$ for
$f_1,f_2\in\hN$, the phases cancelling, so $\ip{\mathcal Rf_1}{\mathcal Rf_2}=\ip{f_1}{f_2}$ by the first paragraph. Thus
$\mathcal R$ is an isometry of $\hN$ into $\mathfrak h'_N$, and it is onto because both spaces have dimension
$L_xL_y/(2\pi)=\Nphi$ (Lemma~\ref{lem:orbitals}(ii)). The operator $\Ff(\mathcal R\otimes\one)$ is then defined as in
Definition~\ref{def:fock}: it maps $c(f_1)^*\cdots c(f_n)^*\Omega_0$ to $c'(f'_1)^*\cdots c'(f'_n)^*\Omega'_0$, with
$f'_i:=(\mathcal R\otimes\one)f_i$ and $\Omega_0$, $\Omega'_0$ the vacua, and it is unitary by the anticommutation
relations, because $\mathcal R\otimes\one$ is unitary.

\emph{The Hamiltonian.} For $f_1,f_2,g_1,g_2\in\hN$ let
\[
\mathcal I(f_1,f_2;g_1,g_2):=4\pi\int_{\Tor_N}(\overline{f_1}g_1)(\overline{f_2}g_2)\,\dd r
+4\pi\int_{\Tor_N}(\overline{f_2}g_2)\,\nabla^2(\overline{f_1}g_1)\,\dd r,
\]
the integrands being smooth functions on $\Tor_N$ by \eqref{eq:bc}, and let $\mathcal I'$ be defined in the same way on
$\Tor'_N$. Then $V_{abcd}=\mathcal I(\psi_a,\psi_b;\psi_c,\psi_d)$ in \eqref{eq:Vabcd}, and by the first paragraph
$\mathcal I'(\mathcal Rf_1,\mathcal Rf_2;\mathcal Rg_1,\mathcal Rg_2)=\mathcal I(f_1,f_2;g_1,g_2)$. For an orthonormal basis
$e=(e_a)$ of $\hN$ let
\[
H[e]:=\frac12\sum_{\sigma,\sigma'}\sum_{a,b,c,d}\mathcal I(e_a,e_b;e_c,e_d)\,c(e_a\otimes\sigma)^*c(e_b\otimes\sigma')^*
c(e_d\otimes\sigma')c(e_c\otimes\sigma).
\]
If $e_a=\sum_ju_{ja}\psi_j$ with a unitary matrix $u$, then, since $\mathcal I$ is conjugate-linear in its first two
arguments and linear in the last two, $c(\cdot)^*$ is linear and $c(\cdot)$ is conjugate-linear,
\begin{multline*}
H[e]=\frac12\sum_{\sigma,\sigma'}\ \sum_{j,k,l,m}\ \sum_{j',k',l',m'}\Bigl(\sum_{a,b,c,d}\overline{u_{ja}}u_{j'a}\,
\overline{u_{kb}}u_{k'b}\,u_{lc}\overline{u_{l'c}}\,u_{md}\overline{u_{m'd}}\Bigr)\\
\times\mathcal I(\psi_j,\psi_k;\psi_l,\psi_m)\,c^*_{j'\sigma}c^*_{k'\sigma'}c_{m'\sigma'}c_{l'\sigma},
\end{multline*}
and the coefficient in parentheses is $\delta_{jj'}\delta_{kk'}\delta_{ll'}\delta_{mm'}$ because $u$ is unitary. So
$H[e]=H[\psi]=H_N$ by \eqref{eq:HN}. In the same way $H'[e']=H'_N$ for every orthonormal basis $e'$ of
$\mathfrak h'_N$, with $H'[e']$ built from $\mathcal I'$ and $c'$. Now $(\mathcal R\psi_a)_a$ is an orthonormal basis of
$\mathfrak h'_N$, conjugation by $\Ff(\mathcal R\otimes\one)$ replaces each $c(\psi_a\otimes\sigma)$ in \eqref{eq:HN} by
$c'(\mathcal R\psi_a\otimes\sigma)$, and $V_{abcd}=\mathcal I(\psi_a,\psi_b;\psi_c,\psi_d)=\mathcal I'(\mathcal R\psi_a,\mathcal R\psi_b;\mathcal R\psi_c,\mathcal R\psi_d)$;
hence
\[
\Ff(\mathcal R\otimes\one)\,H_N\,\Ff(\mathcal R\otimes\one)^*=H'[\mathcal R\psi]=H'_N .
\]
In the same way the number operators, $N_\sigma=\sum_ac(e_a\otimes\sigma)^*c(e_a\otimes\sigma)$ for every orthonormal basis,
are mapped to $N'_\sigma$. So the unitary maps each sector $\sect\nn$ onto the sector of $\mathcal H'_N$ with the same
label and intertwines the restrictions of $H_N$ and $H'_N$, which therefore have the same eigenvalues with
multiplicities.
\end{proof}

\begin{corollary}[uniformity over aspect ratios]\label{cor:uniform}
For even $N\ge L_0^2/(4\pi)$ let $\mathcal T_N$ be the set of rectangular tori with $L_xL_y=4\pi N$ and
$\min\{L_x,L_y\}\ge L_0$, and for $\Tor\in\mathcal T_N$ let $E_{11}(N;\Tor)$ be the eleventh eigenvalue of the hollow-core
Hamiltonian of $\Tor$ in the sector $(N/2,N/2)$. Then $\sup_{\Tor\in\mathcal T_N}E_{11}(N;\Tor)\to0$ as $N\to\infty$ through
the even integers.
\end{corollary}

\begin{proof}
By Lemma~\ref{lem:exchange}, $E_{11}(N;\Tor)$ does not change when the two sides of $\Tor$ are exchanged, so it suffices to
consider the tori of $\mathcal T_N$ with $L_x\ge L_y$. Suppose the claim fails. Then there are $\gamma>0$, even integers
$N_1<N_2<\cdots$ and tori $\Tor_{N_i}\in\mathcal T_{N_i}$ with $L_x\ge L_y$ and $E_{11}(N_i;\Tor_{N_i})\ge\gamma$. Since
$L_x\ge L_y$ and $L_xL_y=4\pi N$, we have $L_x\ge(4\pi N)^{1/2}$; so $\Tor_{N_i}$ is admissible as soon as
$(4\pi N_i)^{1/2}\ge4\delta_0+4$, and we discard the finitely many $i$ for which this fails. Let $\mathcal N$ be the set of
the remaining $N_i$. Then $(\Tor_N)_{N\in\mathcal N}$ are admissible tori with $L_x\ge c_*\sqrt N$ for $c_*:=(4\pi)^{1/2}$,
and the hollow-core Hamiltonians on them satisfy (H1) and (H2) by Corollary~\ref{cor:hc-in-class}, with $C_h=7/\pi$ and
constants $C^H_k$ that depend only on $k$, $\eta$ and $\delta_0$ (Proposition~\ref{prop:slotH}); so
Setting~\ref{set:index} holds. By \eqref{eq:E11} on each of these tori, Assumption~\ref{ass:gap} holds for the sectors
$(N/2,N/2)$, $N\in\mathcal N$, with $p=10$, $\Delta=0$ and $\gamma$. Corollary~\ref{cor:classwide} gives $10\in4\Z$, which is
false.
\end{proof}

In this proof the aspect ratios $L_x/L_y$ of the tori $\Tor_N$, $N\in\mathcal N$, are arbitrary numbers in $[1,\infty)$
and may be unbounded. This is allowed: Setting~\ref{set:index} asks only that the tori be admissible with
$L_x\ge c_*\sqrt N$, and Theorem~\ref{thm:index} and Corollaries~\ref{cor:translation} and \ref{cor:classwide} are proved
under exactly these hypotheses. They hold in this generality because every constant of their proofs depends on the
torus only through admissibility and $c_*$ (Remark~\ref{rem:constants}). The constants of Sections~\ref{sec:frame} and
\ref{sec:hamiltonian} depend only on $\eta$, $\delta_0$ and the order of decay ($\Cvol$ and $C_{\rm str}$ in
Proposition~\ref{prop:geometry}, the constants $C_k$ of Lemma~\ref{lem:localization}, and $C_h$, $C^H_k$ in
Proposition~\ref{prop:slotH}); those of Section~\ref{sec:locality} depend only on these, on the data of the filters, and,
in Lemma~\ref{lem:bookkeeping}, on the separation that is assumed there; and in Sections~\ref{sec:index} and
\ref{sec:translation} the torus enters only through $\abs{\Tor_N}=4\pi N$, $L_y\le\min\{\pi N,(4\pi/c_*)\sqrt N\}$,
$\kappa\le\delta_x\le2\delta_0$, $\ny\le CL_y$ and the separations of Lemma~\ref{lem:separations}, as recorded in the
paragraph on the constants at the end of the proof of Theorem~\ref{thm:index}.

\begin{remark}[thin tori]\label{rem:thin}
Corollary~\ref{cor:uniform} says nothing about tori with $\min\{L_x,L_y\}<L_0$. By Lemma~\ref{lem:exchange} one may take
$L_y$ to be the shorter side; then the slot frame of Section~\ref{sec:frame} is not available, because a window of
width of order $\delta_0$ in the transverse direction need not contain the centre of any orbital
(Lemma~\ref{lem:frame-consistency}(ii) uses $L_y\ge3\pi/\delta_0$), and exchanging the sides instead violates
$L_x\ge4\delta_0+4$ and $L_x\ge c_*\sqrt N$. In particular the index argument proves nothing for sequences of tori
whose circumference $L_y$ stays below $L_0$ while $N\to\infty$. For these tori a variational argument, which uses neither
the slot frame nor the index theorem, bounds $E_{11}(N)$ by an explicit $C(L_y)/N$ (Theorem~\ref{thm:twist}); with
Corollary~\ref{cor:uniform} this gives $E_{11}(N)\to0$ uniformly over all rectangular tori (Corollary~\ref{cor:alltori}).
Section~\ref{subsec:interpretation} comments on this regime.
\end{remark}

\begin{remark}[oblique tori]\label{rem:oblique}
The slot frame, and with it Theorems~\ref{thm:index} and \ref{thm:main}, is constructed here for rectangular tori
only, although Theorem~\ref{thm:kernel} holds for every modular parameter (Appendix~\ref{app:kernel}). For an oblique torus
the Landau-gauge orbitals carry phases quadratic in the orbital index; within a window these split into a linear
part, which shifts the centres of the slots along the cuts and produces a sheared slot lattice, and a remainder
that varies slowly on the scale of a window. We expect the localization estimates of Lemma~\ref{lem:localization}
to hold for the resulting frame, with constants depending on the modular parameter, but this is not
established here.
\end{remark}

\subsection{What is not excluded}

\begin{remark}[a larger gapped ground-state space]\label{rem:loophole}
Theorem~\ref{thm:main} does not exclude the following situation: there are $p\in4\Z$ with $p\ge12$, $\Delta\ge0$ and
$\gamma>2\Delta$ such that, for infinitely many (or for all large) even $N$, $E_p(N)-E_1(N)\le\Delta$ and
$E_{p+1}(N)-E_p(N)\ge\gamma$. This is a uniformly gapped patch of dimension $p$ containing the ten zero modes and
$p-10$ further states of energy at most $\Delta$, with a gap above it; in such a phase the convergence
$E_{11}(N)\to0$ of Theorem~\ref{thm:main}(ii) would be a splitting inside the ground-state space, not the closing of a
gap above it. We do not know whether such a phase exists.

The method of this paper cannot exclude it. With a magnetic translation as the transport unitary, the conclusion of
Theorem~\ref{thm:index} is that $p\,r_N$ lies within $\Onf$ of an integer, where $r_N:=\expP{T_-}$ is the transported
charge per state, which translation covariance determines exactly, as in Lemma~\ref{lem:theta}(d). When $r_N=r$ does
not depend on $N$, the conclusion is $p\,r\in\Z$, and this holds for every multiple $p$ of the denominator of the
rational number $r$. The choices available here give the following values in the sector $(N/2,N/2)$. The up-spin band
charge with $U=\Theta$ gives $r=-n_{\up,N}/\Nphi=-\frac14$ (Corollary~\ref{cor:classwide}). The proofs of
Sections~\ref{sec:hamiltonian}--\ref{sec:translation} use only that the charge is a sum over the orbitals $j$ of the band
of mutually commuting, self-adjoint, sector-preserving operators $n_j$ with integer spectrum, norms bounded
independently of $N$ and $\iota(n_j)\in\Alg(\Col(j))$, whose sum over all $j\in\Z_{\Nphi}$ commutes with every local
term $h_R$, and which satisfy $\Theta^*n_j\Theta=n_{j+1}$; they
therefore apply to the down-spin band charge $\sum_{j\in\Gam}n_{j\dn}$, which gives $r=-n_{\dn,N}/\Nphi=-\frac14$,
and to the total band charge $\sum_{j\in\Gam}(n_{j\up}+n_{j\dn})$, which gives $r=-(n_{\up,N}+n_{\dn,N})/\Nphi=-\frac12$. The
up-spin charge with $U=\Theta^2$, for which $T_-=-n_{0\up}-n_{1\up}$, gives $r=-\frac12$ as well. All of these conditions are
met by $p=12$. Excluding the situation above requires different information: for instance, a lower bound on $E_{11}(N)$
that excludes a splitting inside a gapped ground-state space, or an identification of every uniformly gapped patch at
the bottom of the sector with the space of zero modes.
\end{remark}

\subsection{Interpretation}\label{subsec:interpretation}

\paragraph{What is excluded.}
Theorem~\ref{thm:main} excludes a gapped phase whose ground-state space on the torus is the space of zero modes of the
hollow-core Hamiltonian: on rectangular tori of any fixed aspect ratio, and uniformly over aspect ratios as in
Corollary~\ref{cor:uniform}, the ten zero modes in the sector $(N/2,N/2)$ are not separated from the rest of the
spectrum of this sector by a gap that is uniform in $N$. This is the scenario in which the count of ten would be the
topological ground-state degeneracy of a gapped Haldane--Rezayi phase. The count was first published, from explicit
torus wave functions, by Keski-Vakkuri and Wen \cite{KVW93}, who took it as evidence that the state is non-Abelian;
Gurarie, Flohr and Nayak \cite{GFN97} described the Haldane--Rezayi state by the conformal field theory of central
charge $c=-2$ and attributed the ten-fold degeneracy to a logarithmic partner of the identity field, assuming that the
properties of the zero-energy states of the hollow-core Hamiltonian are stable; and Ino \cite{Ino99} obtained, within
the same description, $5p$ bulk sectors and a torus degeneracy $5p$ at filling $1/p$, hence ten at filling one half, an
argument that presupposes an incompressible state. Read and Rezayi \cite[Section~V]{RR96}, who wrote down the ten
states, held instead that the underlying theory has eight sectors at filling one half, the count of ten arising because
some sectors carry two conformal blocks.
None of these works establishes a gap; Theorem~\ref{thm:main} shows that, for the hollow-core Hamiltonian, the gap above
the ten states that the count of ten as a degeneracy presupposes does not exist.

Nor are the gapped phases with eight states on the torus that have been proposed near the Haldane--Rezayi point
ground-state spaces of $H_N$: the weak- and strong-pairing phases of spin-singlet $d$-wave pairing in the picture of
Read and Green, each with eight states on the torus at filling one half \cite{RG00}, and the unitary description with
eight states of \cite{CGT99}. A uniformly gapped patch at the bottom of the sector contains the ten zero modes, so its
dimension is at least ten and, by Theorem~\ref{thm:main}(i), at least twelve. These proposals satisfy the divisibility
condition of Theorem~\ref{thm:main}(i), but their dimension is too small; the count ten is the one that violates the
condition.

\paragraph{What is not excluded.}
Theorem~\ref{thm:main} does not exclude a gapped phase whose ground-state space strictly contains the zero modes and has
dimension $12$, $16$, $\dots$ (Remark~\ref{rem:loophole}); in such a phase $E_{11}(N)\to0$ would be a splitting inside the
ground-state space. The conclusions that the index argument yields here are divisibility conditions, and such
dimensions meet them. The
expectation since Read and Green \cite{RG00} is stronger than Theorem~\ref{thm:main}: that the hollow-core Hamiltonian has
no uniformly gapped patch at the bottom of the sector at all. Read and Green argued, in a mean-field description of
spin-singlet $d$-wave pairing, that the Haldane--Rezayi state describes the critical point between the weak- and
strong-pairing phases, with bulk fermion excitations whose energy vanishes at zero momentum; their count of states on
the torus at this critical point, five per centre of mass at even $N$ and one spin-$\frac12$ doublet per centre of mass at
odd $N$, is the count of Theorem~\ref{thm:kernel}. Read \cite{Rea09a,Rea09b} gave general arguments that trial states
built from non-unitary conformal field theories do not describe gapped phases, and named the Haldane--Rezayi state as
critical. Exact diagonalization of small systems on the sphere is consistent with a vanishing neutral gap
\cite{CRE19,NPBG23}, and in the limit of a thin torus Seidel and Yang \cite{SY11} constructed excitations whose energy
tends to zero. None of this is proved; Theorem~\ref{thm:main} establishes the part of the picture that concerns the ten
zero modes.

\paragraph{The statement is asymptotic.}
Theorem~\ref{thm:main} concerns the limit $N\to\infty$ only. The constants in Theorem~\ref{thm:index} exist but are not
explicit, and through the light-cone constant $C_{\rm lc}=\max\{1,(4C_1)^4\}$ of Theorem~\ref{thm:propagation}, with $C_1$ bounded by
about $4.5\cdot10^3\,\Cvol\max\{C_h,4C^H_5\}$ (Remark~\ref{rem:propagation}), they are astronomically large. The proof gives neither a
rate for $E_{11}(N)\to0$ nor a value $N_0$ beyond which $E_{11}(N)$ is below a prescribed number, and it says nothing
about any fixed system size, in particular about the sizes accessible to numerical diagonalization.

\paragraph{Thin tori.}
Tori whose circumference $L_y$ stays below $L_0$ while $N\to\infty$ lie outside Corollary~\ref{cor:uniform}
(Remark~\ref{rem:thin}). For them, and more generally at every circumference, Appendix~\ref{app:twist} proves a
variational bound: a zero mode in which the orbitals of the electrons of one spin species are multiplied by a phase that
winds once around the torus gives, for a suitable choice of the zero mode, a state in the sector $(N/2,N/2)$ orthogonal
to all zero modes with energy at most of order $L_y^3/N$ (for $L_y$ bounded below). Hence, with $\kappa=2\pi/L_y$,
\[
E_{11}(N)\le0.40\,(1+0.42\,\kappa)^2\,\frac{L_y^3}N+0.032\,(1+0.41\,\kappa)^2\,\frac{L_y^5}{N^3}
\]
for every even $N$ and every rectangular torus, and the same bound holds for $E_3(N)$ at odd $N$
(Theorem~\ref{thm:twist}). At fixed aspect ratio $L_y^3/N$ grows like $N^{1/2}$, and this bound says nothing; the two
results overlap when $L_0\le L_y\le C$, and the variational bound is needed only for $L_y<L_0$. Together with
Lemma~\ref{lem:exchange} and Corollary~\ref{cor:uniform} it gives $E_{11}(N)\to0$ uniformly over all rectangular tori
(Corollary~\ref{cor:alltori}). Unlike Theorem~\ref{thm:main}, the variational bound is explicit and gives a rate, but only
at bounded or slowly growing circumference; like Theorem~\ref{thm:main}, it does not bear on the situation of
Remark~\ref{rem:loophole}.

\paragraph{The class.}
Corollary~\ref{cor:classwide} is not specific to the hollow-core interaction. For every sequence of Hamiltonians
satisfying (H1) and (H2), the dimension of a uniformly gapped patch at the bottom of the sector $(N/2,N/2)$, in the
sense of Section~\ref{sec:intro} (Assumption~\ref{ass:gap} for these sectors), is divisible by four: this is a fractional Lieb--Schultz--Mattis condition for the up-spin filling $\frac14$ per orbital, in
the sense of \cite[Corollary~5.2]{BBDF21}, and it uses neither frustration-freeness nor zero modes
(Remark~\ref{rem:class-scope}). In particular no Hamiltonian of the class has, in this sector, a uniformly gapped
ground-state space of dimension ten, the dimension that the conformal-field-theory description assigns to a gapped
Haldane--Rezayi phase. Which lowest-Landau-level Hamiltonians other than $H_N$ belong to the class is not examined here.

\paragraph{What remains.}
(a) \emph{A rate at fixed aspect ratio.} Take $p=10$, $\Delta=0$ and $\gamma=E_{11}(N)$. The estimates of Sections~\ref{sec:locality} and
\ref{sec:index}, made explicit as functions of $N$ and $\gamma$ at a single $N$, would bound
$\abs{e^{2\pi\iu\Tr(P\Tt)}-1}$ by a quantity $\varepsilon(N,\gamma)$ that tends to zero as $N\to\infty$ for each $\gamma>0$.
Since $\Tr(P\Tt)=\Tr(PT_-)=-\frac{10}4$ exactly (Lemma~\ref{lem:theta}(d) and (e)), the left-hand side equals $2$, so
$\varepsilon(N,E_{11}(N))\ge2$, an explicit upper bound on $E_{11}(N)$ once $\varepsilon$ is explicit and decreasing in $\gamma$.
A rate obtained in this way would not exclude the situation of Remark~\ref{rem:loophole}. (b) \emph{Explicit constants},
without which no statement about a fixed system size follows. (c) \emph{Oblique tori} (Remark~\ref{rem:oblique}).
(d) \emph{The stronger statement} that $H_N$ has no uniformly gapped patch at the bottom of the sector $(N/2,N/2)$ of any
dimension; it requires excluding the situation of Remark~\ref{rem:loophole}, for which the index argument alone does
not suffice. (e) \emph{Other members of the class} of Definition~\ref{def:class}.
